% Generated draft: requires mathematical and source review.
\documentclass[11pt]{article}
\usepackage[margin=25mm]{geometry}
\usepackage{fontspec}
\setmainfont{Latin Modern Roman}
\usepackage{amsmath,amssymb,mathtools,unicode-math}
\setmathfont{Latin Modern Math}
\usepackage{enumitem,needspace,xurl,hyperref}
\setlength{\parindent}{0pt}
\setlength{\parskip}{5pt}
\setlength{\emergencystretch}{3em}
\newcommand{\N}{\mathbb{N}}
\newcommand{\Z}{\mathbb{Z}}
\newcommand{\Q}{\mathbb{Q}}
\newcommand{\R}{\mathbb{R}}
\newcommand{\C}{\mathbb{C}}
\newcommand{\F}{\mathbb{F}}
\newcommand{\A}{\mathbb{A}}
\newcommand{\PP}{\mathbb{P}}
\newcommand{\E}{\mathbb{E}}
\newcommand{\eps}{\varepsilon}
\newcommand{\dd}{\,\mathrm{d}}
\begin{document}
\section*{UnsolvedMath ID: 30006404}
\textbf{Source problem code:} catalogue entry
\subsection{Definitions and mathematical statement}
The supplied target is to determine the maximum possible size of an \(r\)-cover-free family. The record does not define this term or specify a ground set, its cardinality, or any restriction on the members of the family. Consequently, it does not specify a unique extremal quantity.

A conditional formulation under a conventional finite-set interpretation is the following. Introduce auxiliary integers \(n,r\geq 1\), put \([n]=\{1,\ldots,n\}\), and write \(2^{[n]}\) for its power set. A family \(\mathcal F\subseteq 2^{[n]}\) is called \(r\)-cover-free in this interpretation if, for every choice of pairwise distinct members \(F_0,F_1,\ldots,F_r\in\mathcal F\),
\[
 F_0\not\subseteq\bigcup_{i=1}^{r}F_i.
\]
Define
\[
 C(n,r)=\max\bigl\{|\mathcal F|:\mathcal F\subseteq 2^{[n]}\text{ satisfies this condition}\bigr\}.
\]
The corresponding conditional problem is to determine \(C(n,r)\) for the intended parameter range. The ground-set parameter \(n\), the displayed definition, and this parameter range are formalization choices requiring confirmation from the source; they are not supplied hypotheses. In particular, the displayed condition quantifies over exactly \(r\) other members and is vacuous when \(|\mathcal F|\leq r\). Requiring noncontainment in a union of at most \(r\) other members would require a separately specified convention.

The original excerpt also contains a different conjectural target, attributed to Lovett as Conjecture 5 of reference [11]. It concerns an \(n\times n\) real matrix \(M\) with \(\operatorname{rank}(M)=r\), apparently with at most \(\varepsilon n^2\) nonzero entries, and asserts the existence of an all-zero square submatrix with a stated exponential lower bound on its size. The permitted interval for \(\varepsilon\) and the exponential bound contain corrupted square-root notation, so their exact formulas cannot be stated reliably from this record. Here \(r\) denotes matrix rank; the record does not specify a correspondence with the cover-free parameter above.

The excerpt further defines the intersection matrix of set systems \(\mathcal A,\mathcal B\subseteq 2^{[r]}\) by
\[
 M(A,B)=|A\cap B|\qquad(A\in\mathcal A,\ B\in\mathcal B),
\]
and states that its rank is at most \(r\). Its construction uses the family of all \(k\)-element subsets of \([r]\) and that family's intersection matrix with itself. The formulas selecting \(k\), specifying the matrix order, and describing the claimed matching bound are corrupted. This contextual conjecture and construction cannot supply the missing cover-free hypotheses without further evidence.
\subsection{Short English statement}
Determine the largest size of an \(r\)-cover-free family after specifying the ground set and the precise cover-free convention. A conventional finite-ground-set formulation is provided conditionally. The supplied excerpt also discusses a separate matrix conjecture with damaged formulas, so the intended target and missing hypotheses require human review.
\subsection{Sources}
\begin{itemize}
\item Local UnsolvedMath registry, \texttt{data/unsolvedmath/problems.json}, permanent ID \(30006404\), title \emph{Maximum Sizes of Cover-Free Families}, exact source code \texttt{catalogue entry}. The supplied record lists proposed year 2025 and status \texttt{open}; these are registry attributes, not independently verified mathematical facts. Its \texttt{corrected\_verified} statement label and accompanying typography-repair description are also supplied curation metadata.
\item Supplied report citation: \emph{Combinatorics, Probability and Computing (2026)}, Oberwolfach Reports, DOI 10.4171/owr/2025/42, pp.\ 2243--2322. \url{https://doi.org/10.4171/owr/2025/42}. The record describes Oberwolfach Reports as documenting mathematical workshops and their open-problem discussions. The citation's 2026 label, the DOI's 2025 component, and the proposed year are retained as supplied.
\item Supplied literature source: \emph{On cover-free families of finite vector spaces}. \url{https://doi.org/10.1016/j.jcta.2025.106131}.
\item Supplied literature source: \emph{Oberwolfach Report 42/2025}. \url{https://doi.org/10.4171/owr/2025/42}.
\item Literature assessment checked at \texttt{2026-08-22}: supplied dated curation, not independently verified. It labels the problem open and states that the exact general maximum remains open while recent work determines special classes such as finite vector spaces. The background explicitly calls this dated literature triage derived from the supplied OWR curation and requests independent verification of time-sensitive status and cited claims.
\item The original excerpt attributes its matrix conjecture to Lovett and reference [11], but supplies no complete bibliographic identification for [11]. No proposer is identified for the cover-free-family target. The supplied references have not been visited or independently checked for this draft.
\end{itemize}
\end{document}
